\section{Fixed smoothing and a changing-norm control (NS-63)}
\label{sec:ns63-changing-norm}

\paragraph{Why this check is needed.}
NS-61 uses one fixed extra Mellin integration. Increasing the smoothing
order with the approximation size changes the problem. The following
control shows that relative errors can then vanish with just one fixed
dilation and a fixed coefficient, unconditionally. It is not convergence
in a fixed NB norm.

For an integer $r\geq0$ put
\[
 \tau_r=J^r\chi=\frac{(-\log x)^r}{r!}\mathbf1_{(0,1)},\quad
 \sigma_{n,r}=J^r\rho_n,\quad
 T_r=\norm{\tau_r}^2=\binom{2r}{r}.
\]
Let $E_{N,r}$ be the squared distance from $\tau_r$ to the span of
$\sigma_{1,r},\ldots,\sigma_{N,r}$, so $E_{N,1}=E_N^{(2)}$.

\begin{proposition}[Keep the smoothing order fixed in the RH criterion]
\label{prop:ns63-fixed-order}
For each fixed $r$, $E_{N,r}\to0$ if and only if RH. If
$z=\beta+i\gamma_*$ is a zero with $1/2<\beta<1$, then
\[
 E_{N,r}\geq
 \frac{|z|^{-2(r+1)}}{(2\beta-1)^{-1}+|1-z|^{-2}}.
\]
The lower obstruction itself depends on $r$. After division by $T_r$
it tends to zero for every such $z$, since $|z|>1/2$; it gives no
fixed positive obstruction for the relative error as $r\to\infty$.
\end{proposition}
\begin{proof}
Boundedness of $J^r$ gives $E_{N,r}\leq4^rE_N$ under the original
strong NB criterion. Every $\sigma_{n,r}$ has the same tail $1/(nx)$
for $x>1$, whereas $\tau_r$ has zero tail. Thus the bounded functional
$L_z$ in Proposition~\ref{prop:ns61-rh-equivalence} applies unchanged.
Its value on the target is $z^{-(r+1)}$ and it annihilates the finite
span. This proves the lower bound and the converse for fixed $r$.
\end{proof}

\begin{proposition}[One fixed atom has vanishing relative error as the norm changes]
\label{prop:ns63-one-atom}
Let $z_0=\zeta(1/2)$, $c_0=-1/z_0$, and
\[
 D_0=\frac{\zeta'(1/2)}{\zeta(1/2)}
     =\frac12\left(\gamma+\frac\pi2+\log(8\pi)\right)>0.
\]
Unconditionally, as $r\to\infty$ through integers,
\begin{equation}
 \frac{\norm{\tau_r-c_0\sigma_{1,r}}^2}{T_r}
    =\frac{D_0^2}{4(2r-1)}+O(r^{-2})
    \sim\frac{D_0^2}{8r}\longrightarrow0.
 \label{eq:ns63-relative-control}
\end{equation}
Nevertheless the absolute squared error for this fixed coefficient
diverges:
\begin{equation}
 \norm{\tau_r-c_0\sigma_{1,r}}^2
      \sim\frac{D_0^2}{8\sqrt\pi}\frac{4^r}{r^{3/2}}
      \longrightarrow\infty.
 \label{eq:ns63-absolute-control}
\end{equation}
Moreover $(J/2)^rf\to0$ in $H$ for every fixed $f\in H$.
Neither this normalized convergence nor
\eqref{eq:ns63-relative-control} tests RH.
\end{proposition}
\begin{proof}
The alternating-series expression for zeta shows $z_0<0$: its
numerator $\sum_{n\geq1}(-1)^{n-1}n^{-1/2}$ is positive and its
denominator $1-\sqrt2$ is negative. Thus $c_0$ is well-defined.
Mellin Plancherel writes the relative squared error as the average of
\[
 H(t)=\left|1-\frac{\zeta(1/2+it)}{z_0}\right|^2
\]
under the probability density proportional to
$(1/4+t^2)^{-(r+1)}dt$. The exact second and fourth moments of this
density are
\[
 \mathbb E_r(t^2)=\frac1{4(2r-1)},\qquad
 \mathbb E_r(t^4)=\frac3{16(2r-1)(2r-3)}
\]
for $r\geq1$ and $r\geq2$, respectively. They follow directly
from the beta integral. Near zero, conjugation symmetry and Taylor
expansion give $H(t)=D_0^2t^2+O(t^4)$.
Away from any fixed neighborhood of zero the normalized integral is
exponentially small in $r$. To justify this without an unproved zeta
estimate, use
\[
 \left|\frac{\zeta(1/2+it)}{1/2+it}\right|
 \leq\int_0^\infty\rho_1(x)x^{-1/2}\,dx<4,
\]
or the integrability of the original Mellin square; either controls
the entire exterior integral. The displayed moments then prove
\eqref{eq:ns63-relative-control}. Differentiating the zeta functional
equation at $1/2$, and using
$\psi(1/2)=-\gamma-2\log2$ \cite[Eq. (5.4.13)]{DLMFGammaBounds},
gives the stated value of $D_0$.
Finally $T_r=\binom{2r}{r}\sim4^r/\sqrt{\pi r}$ proves
\eqref{eq:ns63-absolute-control}.

For the last assertion the Mellin multiplier of $(J/2)^r$ is
$(2s)^{-r}$. Its modulus is at most one on $\Re s=1/2$ and tends
to zero for every $t\ne0$. Dominated convergence in the Mellin
$L^2$ norm applies to every fixed $f$. The exceptional point $t=0$
has measure zero.
\end{proof}

\paragraph{Finite control and its full tail.}
The companion Arb calculation integrates the analytic product
\[
 \frac{(1-\zeta(s)/z_0)(1-\zeta(1-s)/z_0)}
             {\pi T_r\{s(1-s)\}^{r+1}},\qquad s=\tfrac12+it,
\]
over $0\leq t\leq A$. It uses reflection $1-s$, not complex
conjugation of the integration variable, so the quadrature callback
is meromorphic. The preceding elementary bound supplies the complete
positive tail bound
\begin{equation}
 \frac1{\pi T_r}\left\{
       \frac{2A^{-2r-1}}{2r+1}
       +\frac{32A^{-2r+1}}{z_0^2(2r-1)}\right\},\qquad r\geq1.
 \label{eq:ns63-quadrature-tail}
\end{equation}
The fixed-order interval results illustrate the control, not the
asymptotic proof or any RH assertion.

The useful research problem is therefore the fixed-order cofinal
correlation estimate in~\eqref{eq:ns61-open-correlation}. Further
smoothing cannot be credited as a convergence result without retaining
one fixed norm and its target. The unsolved arithmetic estimate remains
the one for the optimized coefficients as $N$ increases.
